\chapter{Fiscalité minière et modèles d'analyse fiscale~: cadre conceptuel}
\label{chap:cadre}

\section{La fiscalité minière dans les économies riches en ressources}

La conception d'un régime fiscal applicable aux industries extractives répond à un arbitrage récurrent dans les pays en développement dotés de ressources naturelles~: maximiser la part de la rente économique captée par la puissance publique sans décourager des investissements dont le montant, l'horizon et le risque sont sensiblement plus élevés que dans la plupart des autres secteurs \citep{charlet2013}. Un régime trop favorable à l'investisseur prive l'État de recettes qu'il aurait pu légitimement percevoir~; un régime trop contraignant peut conduire au report, voire à l'abandon, de projets par ailleurs viables d'un point de vue géologique et technique \citep{otto2006royalties}.

Cet arbitrage est d'autant plus sensible que le secteur extractif présente des caractéristiques particulières~: une forte intensité capitalistique concentrée dans les premières années du projet, une volatilité importante des prix des matières premières, une asymétrie d'information entre l'État et l'investisseur quant aux coûts réels du projet, et un horizon temporel qui dépasse largement le cycle électoral ou budgétaire \citep{nrgi2015fiscalregime}. La Charte des ressources naturelles, promue par NRGI, rappelle à cet égard que la production de recettes publiques constitue le principal canal par lequel l'exploitation des ressources naturelles peut contribuer durablement au développement, à condition que le système fiscal retenu prenne en compte la nature spécifique de ces ressources, l'incertitude entourant leur exploitation et les capacités de l'administration fiscale \citep{nrgi2014charter}.

Dans le cas spécifique de l'Afrique de l'Ouest, plusieurs travaux soulignent que la générosité des incitations fiscales accordées aux sociétés extractives --- exonérations temporaires, clauses de stabilité, conventions particulières dérogeant au droit commun --- peut représenter un manque à gagner budgétaire considérable. \citet{actionaid2015giveaway} estiment ainsi à plusieurs milliards de dollars par an, pour le seul groupe formé par le Ghana, le Nigeria et le Sénégal, le coût des incitations fiscales accordées aux sociétés multinationales. Ce constat justifie l'intérêt porté, depuis le milieu des années 2010, aux outils de modélisation permettant de quantifier \emph{ex ante} l'impact d'un régime fiscal donné --- ou d'une convention minière particulière --- sur le partage de la rente entre l'État et l'investisseur.

\section{Typologie des instruments fiscaux applicables aux projets miniers}
\label{sec:typologie}

La littérature distingue classiquement les instruments fiscaux miniers selon deux critères~: leur assiette (volume produit, valeur des ventes, bénéfice comptable, flux de trésorerie, ou capital de la société) et le moment du cycle de vie du projet auquel ils s'appliquent (phase de recherche ou phase d'exploitation) \citep{otto2006royalties}. On retrouve ainsi, parmi les instruments les plus répandus dans les régimes miniers africains~:

\begin{itemize}
  \item les \textbf{redevances superficiaires}, dues à l'octroi et au renouvellement des titres miniers, dont le rendement budgétaire est généralement marginal~;
  \item les \textbf{redevances minières} (\emph{royalties}), assises sur le volume ou la valeur de la production, brute ou nette de certains coûts (transport, traitement)~; elles garantissent à l'État un paiement minimal indépendant de la profitabilité du projet \citep{guj2012royalties}~;
  \item l'\textbf{impôt sur les sociétés} (ou impôt sur les bénéfices industriels et commerciaux), assis sur le résultat fiscal et qui permet, par construction, un partage du risque entre l'État et l'investisseur~;
  \item la \textbf{participation de l'État} au capital de la société d'exploitation, gratuite ou payante, dont le rendement dépend de la distribution effective de dividendes~;
  \item les \textbf{retenues à la source sur les dividendes} (impôt sur le revenu des valeurs mobilières), perçues au moment du rapatriement des bénéfices~;
  \item les \textbf{droits de douane} sur les biens d'équipement, intrants et consommables importés, fréquemment exonérés pour la phase de construction dans le cadre des conventions minières~;
  \item enfin, des instruments dits « progressifs » --- impôts sur la rente, sur les profits excédentaires ou redevances à taux variable --- dont le taux ou l'assiette évolue avec la profitabilité du projet ou le niveau des prix, et qui permettent à l'État de bénéficier d'une part accrue des recettes lorsque la conjoncture des prix est favorable, tout en limitant la charge fiscale en période de prix bas \citep{rota-graziosi2015mali}.
\end{itemize}

Les comparaisons internationales mobilisées dans ce rapport (chapitre~\ref{chap:resultats}) reposent sur des combinaisons différentes de ces instruments. À titre d'illustration, la redevance minière varie de 0\,\% (Pérou, sur option) à 5\,\% (Ghana, Sénégal, Zambie) tandis que le taux nominal de l'impôt sur les bénéfices s'échelonne entre 22,4\,\% (Chili) et 35\,\% (Ghana, RDC amendée) \citep{pwc2012,laporte2016or}. La participation gratuite de l'État, en revanche, n'est observée que dans les juridictions en développement~: 10\,\% au Sénégal et dans la RDC de 2002, contre 10\,\% également en RDC amendée mais combinée à d'autres instruments plus lourds.

\section{Approches d'évaluation de la rentabilité et de la performance fiscale}

L'évaluation \emph{ex ante} de la rentabilité d'un projet minier et de la performance d'un régime fiscal mobilise principalement l'approche des flux de trésorerie actualisés (\emph{Discounted Cash Flow}, DCF), qui projette, sur la durée de vie du projet, les flux de trésorerie avant et après impôt et les actualise à un taux représentatif du coût d'opportunité du capital \citep{otto2006royalties}. Cette approche présente l'avantage de la simplicité et de la transparence~: elle permet de faire apparaître explicitement chaque paramètre technique, économique et fiscal, et de comparer aisément plusieurs scénarios de régime fiscal sur la base d'un même profil opérationnel.

Des approches alternatives ont été développées pour pallier la principale limite du DCF --- l'utilisation d'un taux d'actualisation unique, indépendant de la nature du risque associé à chaque flux. La méthode dite \emph{Modern Asset Pricing} (MAP), forme simplifiée des options réelles (\emph{Real Option Valuation}, ROV), traite séparément le risque de prix, généralement modélisé par un processus stochastique (\emph{Geometric Brownian Motion}, modèles à retour à la moyenne, etc.), et les autres sources de risque \citep{guj2012royalties}. La méthode \emph{Decoupled Net Present Value} (DNPV), plus récente, combine quant à elle les risques de marché et les risques non marchands (climatiques, sismiques, sociaux) pour évaluer les opportunités d'investissement minier. Ces approches, plus exigeantes en données et en complexité de mise en \oe uvre, demeurent peu utilisées par les administrations publiques des pays en développement, en raison notamment de la confidentialité des données nécessaires à leur calibration \citep{kourouma2019}.

C'est pourquoi le FMI privilégie, pour l'assistance technique apportée aux pays riches en ressources, l'approche DCF formalisée dans le modèle FARI, présenté à la section suivante. Le prix de la ressource y est supposé constant en termes réels sur toute la durée de vie du projet --- hypothèse forte mais usuelle, dont la robustesse peut être testée au moyen d'analyses de sensibilité sur le prix et sur les coûts.

\section{Le modèle FARI et son adaptation par NRGI}
\label{sec:fari}

\subsection{Principes généraux}

Le modèle \emph{Fiscal Analysis of Resource Industries} (FARI) est l'outil de référence utilisé par le Département des finances publiques du FMI pour la conception, l'évaluation et la prévision des régimes fiscaux applicables aux industries extractives \citep{luca2016fari}. Il est mobilisé dans une trentaine de pays riches en ressources, à l'occasion de l'ouverture de nouveaux gisements, de la renégociation de conventions existantes, ou encore de l'analyse de l'écart entre les recettes effectivement perçues et celles prévues par le modèle. Sa structure peut être résumée en trois blocs~: des \emph{intrants} (données techniques et économiques du projet, paramètres du régime fiscal, hypothèses financières), des \emph{calculs} (flux de trésorerie avant impôt, calcul successif de chaque instrument fiscal, flux après impôt pour l'investisseur et pour l'État) et des \emph{résultats} (indicateurs standardisés tels que la \VAN, le \TRI après impôt et le \TEMI).

NRGI a publié une version simplifiée de cet outil sous la forme d'un classeur Excel, structurée en quatre feuilles --- présentation, paramètres, résultats et modèle de calcul --- et paramétrée par défaut pour le code minier de la République Démocratique du Congo, avec dix scénarios de comparaison internationale intégrés. C'est cette version, dont l'intégrité méthodologique est préservée par NRGI tant que les adaptations n'altèrent pas la logique de calcul \citep{nrgi_excel_or}, qui constitue le support de la présente étude (voir chapitre~\ref{chap:methodo}).

\subsection{Formalisation}

Pour une année $t$ comprise entre le début et la fin de vie du projet, le flux de trésorerie du projet avant toute imposition s'écrit~:
\begin{equation}
  \mathrm{FT}^{\text{avant impôt}}_t = Q_t \cdot P_t - I_t - C_t - F_t
  \label{eq:fari1}
\end{equation}
où $Q_t$ désigne le volume produit et $P_t$ le prix unitaire du minerai à l'année $t$, $I_t$ les dépenses d'investissement (initiales et de maintien), $C_t$ les coûts d'exploitation et $F_t$ les coûts de fermeture et de réhabilitation du site. Le flux de trésorerie après impôt revenant à l'investisseur s'obtient en déduisant de \eqref{eq:fari1} l'ensemble des prélèvements fiscaux et parafiscaux $T_t$ (redevances, impôts, droits, participation de l'État, etc.)~:
\begin{equation}
  \mathrm{FT}^{\text{après impôt}}_t = \mathrm{FT}^{\text{avant impôt}}_t - T_t
  \label{eq:fari2}
\end{equation}

La valeur actuelle nette des flux de trésorerie de l'investisseur, qui constitue le critère de décision d'investissement, est donnée par~:
\begin{equation}
  \VAN_{\text{investisseur}} = \sum_{t=0}^{n} \frac{\mathrm{FT}^{\text{après impôt}}_t}{(1+r)^t}
  \label{eq:fari3}
\end{equation}
où $r$ est le taux d'actualisation de l'investisseur et $n$ la dernière année du projet. La valeur actuelle nette des recettes fiscales et parafiscales perçues par l'État sur la même période s'écrit~:
\begin{equation}
  \VAN_{\text{État}} = \sum_{t=0}^{n} \frac{T_t}{(1+r')^t}
  \label{eq:fari4}
\end{equation}
avec $r'$ le taux d'actualisation retenu du point de vue de l'État (dans le modèle utilisé, $r'=10\,\%$ et $r=12{,}5\,\%$, conformément aux hypothèses de référence du modèle FARI).

Le \TEMI, indicateur central pour comparer la performance de différents régimes fiscaux entre eux ou avec des régimes étrangers, est défini comme le rapport entre la valeur actuelle nette des recettes de l'État et celle du flux de trésorerie du projet avant impôt~:
\begin{equation}
  \TEMI = \frac{\VAN_{\text{État}}}{\VAN\!\left(\mathrm{FT}^{\text{avant impôt}}\right)}
  \label{eq:fari5}
\end{equation}

Un \TEMI élevé traduit une charge fiscale importante pour l'investisseur~; un \TEMI faible signale, à l'inverse, que l'État retient une part relativement modeste de la rente générée par le projet. \citet{lassourd2018rdc} souligne qu'un \TEMI sensiblement supérieur aux standards internationaux --- de l'ordre de 40 à 50\,\% pour les projets aurifères selon les comparaisons usuelles --- peut freiner le développement de projets par ailleurs économiquement viables, sans pour autant garantir une mobilisation optimale des recettes pour l'État si le régime décourage l'investissement initial.

À ces indicateurs s'ajoutent, du point de vue de l'investisseur, le \TRI après impôt --- comparé au taux de rendement minimal exigé --- et la période de récupération de l'investissement, qui mesurent respectivement la rentabilité intrinsèque du projet et la rapidité avec laquelle les capitaux engagés sont récupérés. Le chapitre~\ref{chap:resultats} reviendra sur les conditions dans lesquelles ces deux derniers indicateurs restent interprétables.
